% Part:sets-functions-relations
% Chapter: size-of-sets
% Section: non-enumerability

\documentclass[../../../include/open-logic-section]{subfiles}

\begin{document}

\olfileid{sfr}{siz}{nen}

\olsection{\printtoken{S}{nonenumerable} verzamelingen}

\begin{editorial}
  Deze paragraaf bewijst met de definitie uit \olref[enm]{sec} de
  overaftelbaarheid van $\Bin^\omega$ en $\Pow{\PosInt}$. De behandeling
  is iets elementairder en iets uitvoeriger opgezet dan de versie in
  \olref[enm-alt]{sec}
\end{editorial}

Sommige verzamelingen, zoals de verzameling $\PosInt$ van positieve
gehele getallen, zijn oneindig. Tot dusver hebben we voorbeelden gezien
van oneindige verzamelingen die allemaal !!{enumerable} waren. Er zijn
echter ook oneindige verzamelingen die deze eigenschap niet hebben. Zulke
verzamelingen heten \emph{!!{nonenumerable}}.

Om te beginnen is het wellicht al verrassend dat er
!!{nonenumerable} verzamelingen bestaan. Voor iedere !!{enumerable}
verzameling~$A$ bestaat er een !!a{surjective} functie $f \colon \PosInt
\to A$. Als een verzameling !!{nonenumerable} is, bestaat zo'n functie
niet. Geen enkele functie die de oneindig vele !!{element}s van~$\PosInt$
op~$A$ afbeeldt, kan dan dus heel~$A$ bereiken. In die zin heeft~$A$
``meer'' !!{element}s dan de oneindig vele positieve gehele getallen.

Hoe kan men bewijzen dat een verzameling !!{nonenumerable} is? Daarvoor
moet men aantonen dat zo'n surjectieve functie niet kan bestaan. Dit is
gelijkwaardig met aantonen dat de elementen van~$A$ niet kunnen worden
opgesomd in een lijst die naar één kant oneindig doorloopt. De beste
aanpak is te bewijzen dat iedere lijst van !!{element}s van~$A$ ten
minste één element weglaat; met andere woorden, dat geen functie
$f\colon \PosInt \to A$ !!{surjective} kan zijn. Dit kan met Cantors
\emph{diagonaalmethode}. Gegeven een lijst van !!{element}s van~$A$,
bijvoorbeeld $x_1$, $x_2$, \dots, construeren we een ander element
van~$A$ dat door zijn constructie onmogelijk op die lijst kan staan.

Ons eerste voorbeeld is de verzameling~$\Bin^\omega$ van alle oneindige
rijen van $0$'en en $1$'en zonder ontbrekende plaatsen.

\begin{thm}
\ollabel{thm:nonenum-bin-omega}
De verzameling $\Bin^\omega$~is !!{nonenumerable}.
\end{thm}

\begin{proof}
Veronderstel ter tegenspraak dat $\Bin^\omega$ !!{enumerable} is. Neem
dus aan dat $s_{1}$, $s_{2}$, $s_{3}$, $s_{4}$, \dots{} een lijst van
alle !!{element}s van~$\Bin^\omega$ is. Iedere $s_i$ is zelf een
oneindige rij van $0$'en en~$1$'en. Noem het $j$-de element van de
$i$-de rij in deze lijst $s_i(j)$. De $i$-de rij~$s_i$ is dan
\[
s_i(1), s_i(2), s_i(3), \dots
\]

We kunnen deze lijst en de elementen van iedere rij $s_i$ daarin in een
rooster rangschikken:
\[
\begin{array}{c|c|c|c|c|c}
& 1 & 2 & 3 & 4 & \dots \\\hline
1 & \mathbf{s_{1}(1)} & s_{1}(2) & s_{1}(3) & s_1(4) & \dots \\\hline
2 & s_{2}(1)& \mathbf{s_{2}(2)} & s_2(3) & s_2(4) & \dots \\\hline
3 & s_{3}(1)& s_{3}(2) & \mathbf{s_3(3)} & s_3(4) & \dots \\\hline
4 & s_{4}(1)& s_{4}(2) & s_4(3) & \mathbf{s_4(4)} & \dots \\\hline
\vdots & \vdots & \vdots & \vdots & \vdots & \mathbf{\ddots}
\end{array}
\]
De aanduidingen links geven het nummer van de rij in de lijst $s_1$,
$s_2$, \dots; de getallen bovenaan nummeren de !!{element}s van de
afzonderlijke rijen. Zo is $s_{1}(1)$ de naam van het getal, een $0$ of
een~$1$, dat het eerste !!{element} van de rij $s_{1}$ is, enzovoort.

Nu construeren we een oneindige rij $\overline{s}$ van $0$'en en $1$'en
die onmogelijk op deze lijst kan staan. De definitie van $\overline{s}$
hangt af van de lijst $s_1$, $s_2$, \dots. Iedere oneindige lijst van
oneindige rijen van $0$'en en $1$'en bepaalt een oneindige
rij~$\overline{s}$ die gegarandeerd niet op de lijst voorkomt.

Om $\overline{s}$ te definiëren, leggen we al haar !!{element}s vast:
we geven $\overline{s}(n)$ voor iedere $n \in \PosInt$. We lezen daartoe
de diagonaal van het bovenstaande rooster af (vandaar de naam
``diagonaalmethode'') en vervangen vervolgens iedere $1$ door een $0$
en iedere $0$ door een~$1$. Abstract geformuleerd definiëren we
$\overline{s}(n)$ als $0$ of $1$, naargelang het $n$-de !!{element} op
de diagonaal, $s_n(n)$, gelijk is aan $1$ of $0$.
\[
\overline{s}(n) =
\begin{cases}
1 & \text{als $s_{n}(n) = 0$}\\
0 & \text{als $s_{n}(n) = 1$}.
\end{cases}
\]
Wie formules verkiest boven definities per geval, kan ook
$\overline{s}(n) = 1 - s_n(n)$ definiëren.

De rij $\overline{s}$ is duidelijk een oneindige rij van $0$'en en
$1$'en: zij ontstaat uit de rij van $0$'en en $1$'en op de diagonaal
van ons rooster door iedere term om te keren. Dus is $\overline{s}$
!!a{element} van~$\Bin^\omega$. Zij kan echter niet op de lijst $s_1$,
$s_2$, \dots{} staan. Waarom niet?

Zij kan niet de eerste rij op de lijst, $s_1$, zijn, want zij verschilt
van $s_1$ in het eerste !!{element}. Wat $s_1(1)$ ook is, we hebben
voor $\overline{s}(1)$ de andere waarde gekozen. Zij kan evenmin de
tweede rij op de lijst zijn, want $\overline{s}$ verschilt van $s_2$ in
het tweede element: als $s_2(2)$ gelijk is aan $0$, is $\overline{s}(2)$
gelijk aan $1$, en omgekeerd. Enzovoort.

Preciezer: als $\overline{s}$ op de lijst stond, zou er een $k$ bestaan
waarvoor $\overline{s} = s_{k}$. Twee rijen zijn precies dan gelijk als
ze op iedere plaats overeenkomen; voor iedere~$n$ zou dus
$\overline{s}(n) = s_{k}(n)$ gelden. In het bijzonder zou voor het
speciale geval $n = k$ moeten gelden dat $\overline{s}(k) = s_{k}(k)$.
Nu is $s_k(k)$ gelijk aan $0$ of aan~$1$. Als die waarde $0$ is, moet
$\overline{s}(k)$ gelijk zijn aan~$1$---zo hebben we $\overline{s}$
gedefinieerd. Maar als $s_k(k) = 1$, dan volgt opnieuw uit de manier
waarop we $\overline{s}$ hebben gedefinieerd dat $\overline{s}(k) = 0$.
In beide gevallen geldt $\overline{s}(k) \neq
s_{k}(k)$.

We begonnen met de aanname dat er een lijst van !!{element}s van
$\Bin^\omega$, $s_1$, $s_2$, \dots{}, bestaat. Uit deze lijst
construeerden we
een rij~$\overline{s}$ waarvan we hebben bewezen dat zij niet op de
lijst kan staan. Zij is echter beslist zelf een rij van $0$'en en
$1$'en als alle $s_i$ rijen van $0$'en en $1$'en zijn; dus $\overline{s} \in
\Bin^\omega$. Er kan in het bijzonder geen lijst van \emph{alle}
!!{element}s van~$\Bin^\omega$ bestaan. Uit iedere vermeende lijst
kunnen we immers ook een rij~$\overline{s}$ construeren die gegarandeerd
niet op die lijst staat. De aanname dat alle rijen in~$\Bin^\omega$ op
één lijst staan, leidt dus tot een tegenspraak.
\end{proof}

\begin{explain}
Deze bewijsmethode heet een ``diagonaalargument'', omdat de diagonaal van
het rooster wordt gebruikt om~$\overline{s}$ te definiëren. Voor een
diagonaalargument is niet noodzakelijk een rooster nodig: met hetzelfde
idee kunnen we bewijzen dat verzamelingen niet !!{enumerable} zijn, ook
als er geen rooster en geen daadwerkelijke diagonaal aan te pas komt.
\end{explain}

\begin{thm}
\ollabel{thm:nonenum-pownat}
$\Pow{\PosInt}$ is niet !!{enumerable}.
\end{thm}

\begin{proof}
We gaan op dezelfde manier te werk en laten zien dat er voor iedere
lijst van deelverzamelingen van~$\PosInt$ een deelverzameling van
$\PosInt$ bestaat die niet op de lijst kan staan. Veronderstel dat de
volgende lijst van deelverzamelingen van~$\PosInt$ gegeven is:
\[
Z_{1}, Z_{2}, Z_{3}, \dots
\]
We definiëren nu een verzameling $\overline{Z}$ waarvoor bij iedere
$n \in \PosInt$ geldt dat $n \in \overline{Z}$ precies dan als $n
\notin Z_{n}$:
\[
\overline{Z} = \Setabs{n \in \PosInt}{n \notin Z_n}
\]
$\overline{Z}$ is duidelijk een verzameling van positieve gehele
getallen, want volgens de aanname is iedere~$Z_n$ dat ook; bijgevolg
geldt $\overline{Z} \in \Pow{\PosInt}$. Maar $\overline{Z}$ kan niet op
de lijst staan. Om dit aan te tonen, bewijzen we voor iedere $k \in
\PosInt$ dat $\overline{Z} \neq Z_k$.

Laat $k \in \PosInt$ willekeurig zijn. We hebben $\overline{Z}$ zo
gedefinieerd dat voor iedere $n \in \PosInt$ geldt: $n \in
\overline{Z}$ precies dan als $n \notin Z_n$. Als speciaal geval met
$n=k$ geldt dus: $k \in \overline{Z}$ precies dan als $k \notin Z_k$.
Daaruit volgt dat $\overline{Z} \neq Z_k$: $k$ is !!a{element} van de
ene verzameling, maar niet van de andere, zodat $\overline{Z}$ en $Z_k$
verschillende !!{element}s hebben. Omdat $k$ willekeurig was, staat
$\overline{Z}$ niet op de lijst $Z_1$, $Z_2$, \dots
\end{proof}

\begin{explain}
In het voorgaande bewijs werd geen diagonaal genoemd, maar men kan het
wel als een diagonaalargument opvatten. Stel daartoe de verzamelingen
$Z_1$, $Z_2$, \dots{} in een rooster geschreven voor, waarbij ieder
!!{element}~$j \in Z_i$ in de $j$-de kolom staat. Stel bijvoorbeeld dat
de eerste vier verzamelingen op de lijst $\{1,2,3,\dots\}$, $\{2, 4, 6,
\dots\}$, $\{1,2,5\}$ en $\{3,4,5,\dots\}$ zijn. Dan begint het rooster
als volgt:
\[
\begin{array}{r@{}rrrrrrr}
  Z_1 = \{ & \mathbf{1}, & 2, & 3, & 4, & 5, & 6, & \dots\}\\
  Z_2 = \{ &  & \mathbf{2}, &  & 4, &  & 6, & \dots\}\\
  Z_3 = \{ & 1, & 2, &  &  & 5\phantom{,} &  & \}\\
  Z_4 = \{ &  &  & 3, & \mathbf{4}, & 5, & 6, & \dots\}\\
  \vdots & & & & & \ddots
\end{array}
\]
De verzameling $\overline{Z}$ wordt dan verkregen door de diagonaal af
te lopen, de getallen op de diagonaal weg te laten en juist die $j$ op
te nemen waarvoor rij en kolom nummer~$j$ elkaar in het rooster op een
lege plaats kruisen. In dit voorbeeld laten we $1$ en $2$ weg, nemen
we~$3$ op, laten we~$4$ weg, enzovoort.
\end{explain}

\begin{prob}
Bewijs met een diagonaalargument dat $\Pow{\Nat}$ !!{nonenumerable} is.
\end{prob}

\begin{prob}\label{sfr:siz:nen:prob:f-posint}
Bewijs met een expliciet diagonaalargument dat de verzameling functies
$f \colon \PosInt \to \PosInt$ !!{nonenumerable} is. Toon dus aan dat er,
als $f_1$, $f_2$, \dots{} een lijst van functies is en iedere $f_i\colon
\PosInt \to \PosInt$ een functie is, een $\overline{f}\colon \PosInt \to \PosInt$
bestaat die niet op deze lijst staat.
\end{prob}
\end{document}
